\section{Risultati sperimentali}
\label{sec:results}

Questa sezione sarà compilata dopo l'esecuzione delle campagne. I pilot locali
servono a verificare le pipeline e a formulare ipotesi; le conclusioni
confermative saranno basate sui risultati ottenuti con i protocolli e le
configurazioni congelati. Ogni tabella deve riportare il numero di run
complete, le FE effettive, i seed e le funzioni incluse; i risultati non
completi non devono essere aggregati come se fossero una campagna conclusa.

\subsection{Equivalenza delle modalità sincrone}
\label{sec:results-equivalence}

Da riportare: percentuale di run appaiate con identità esatta di best value,
posizione e cronologia best-so-far; eventuali mismatch, loro riproducibilità e
causa. L'esito deve essere verificato per master--slave e particle-parallel a
ciascun numero di worker. Non è sufficiente confrontare soltanto il valore
finale.

\IfFileExists{figures/internal_pilots/fig_01_synchronous_equivalence.pdf}{%
\begin{figure}[t]
  \centering
  \includegraphics[width=\columnwidth]{figures/internal_pilots/fig_01_synchronous_equivalence.pdf}
  \caption{Controllo dell'equivalenza esatta tra esecutori sincroni.}
  \label{fig:synchronous-equivalence}
\end{figure}}{\noindent\fbox{\parbox{0.95\columnwidth}{\centering Da generare: \texttt{fig\_01\_synchronous\_equivalence.pdf}.}}}

\subsection{Scaling e costo dell'esecutore}
\label{sec:results-scaling}

Da riportare: tempi appaiati, speedup, efficienza parallela e intervalli di
incertezza per master--slave e particle-parallel. Separare i tre livelli di
costo della funzione obiettivo (nessun lavoro aggiuntivo, carico intermedio e
carico elevato) e le dimensioni/population size previste. I tempi assoluti
sono riferiti esclusivamente alla macchina descritta nel manifest.

\IfFileExists{figures/internal_pilots/fig_02_strong_scaling.pdf}{%
\begin{figure}[t]
  \centering
  \includegraphics[width=\columnwidth]{figures/internal_pilots/fig_02_strong_scaling.pdf}
  \caption{Strong scaling degli esecutori CPSO sincroni.}
  \label{fig:strong-scaling}
\end{figure}}{\noindent\fbox{\parbox{0.95\columnwidth}{\centering Da generare: \texttt{fig\_02\_strong\_scaling.pdf}.}}}

\IfFileExists{figures/internal_pilots/fig_03_parallel_efficiency.pdf}{%
\begin{figure}[t]
  \centering
  \includegraphics[width=\columnwidth]{figures/internal_pilots/fig_03_parallel_efficiency.pdf}
  \caption{Efficienza parallela per i backend sincroni.}
  \label{fig:parallel-efficiency}
\end{figure}}{}

\subsection{Aggiornamento asincrono}
\label{sec:results-async}

Da riportare: qualità best-so-far e successo a budget FE fissato, distribuzione
dei tempi, sensibilità al numero di worker e alla topologia sociale. Il
confronto con il seriale è algoritmico, non uno speedup a traiettoria invariata.

\IfFileExists{figures/internal_pilots/fig_05_asynchronous_quality.pdf}{%
\begin{figure}[t]
  \centering
  \includegraphics[width=\columnwidth]{figures/internal_pilots/fig_05_asynchronous_quality.pdf}
  \caption{Confronto esplorativo della qualità per la modalità event-driven.}
  \label{fig:asynchronous-quality}
\end{figure}}{\noindent\fbox{\parbox{0.95\columnwidth}{\centering Da generare: \texttt{fig\_05\_asynchronous\_quality.pdf}.}}}

\subsection{Architettura a isole}
\label{sec:results-islands}

Da riportare: confronto FE-matched al variare di numero di isole, topologia
intra-isola e politica di migrazione; quota di migrazioni accettate e costi di
comunicazione, se disponibili. Distinguere l'effetto della decomposizione da
quello della migrazione.

\IfFileExists{figures/internal_pilots/fig_06_island_quality.pdf}{%
\begin{figure}[t]
  \centering
  \includegraphics[width=\columnwidth]{figures/internal_pilots/fig_06_island_quality.pdf}
  \caption{Confronto esplorativo della qualità nel modello a isole.}
  \label{fig:island-quality}
\end{figure}}{\noindent\fbox{\parbox{0.95\columnwidth}{\centering Da generare: \texttt{fig\_06\_island\_quality.pdf}.}}}

\subsection{Espansione alle sette topologie}
\label{sec:results-topologies}

Da riportare: qualità per funzione e dimensione, ranghi aggregati con unità
statistiche corrette, costo della selezione sociale e verifica dell'equivalenza
tra esecutori sincroni per topologia. La campagna a sette topologie non è
confusa con la campagna architetturale principale a tre topologie.

\IfFileExists{figures/internal_pilots/fig_08_current_topology_quality.pdf}{%
\begin{figure}[t]
  \centering
  \includegraphics[width=\columnwidth]{figures/internal_pilots/fig_08_current_topology_quality.pdf}
  \caption{Confronto delle sette famiglie topologiche nel pilot dedicato.}
  \label{fig:seven-topologies}
\end{figure}}{\noindent\fbox{\parbox{0.95\columnwidth}{\centering Da generare: \texttt{fig\_08\_current\_topology\_quality.pdf}.}}}

\subsection{Benchmark esterni: COCO BBOB e BBOB-LargeScale}
\label{sec:results-bbob}

Per ogni suite saranno riportate le ECDF e le misure runtime-to-target prodotte
con il post-processing ufficiale COCO. Le tabelle accessorie con i ranghi
finali sono descrittive e non sostituiscono le metriche anytime. I risultati
BBOB standard e BBOB-LargeScale saranno analizzati separatamente.

\IfFileExists{figures/external_pilots/rank_cpso_bbob_stageC_pilot_v1_final_best_value.pdf}{%
\begin{figure}[t]
  \centering
  \includegraphics[width=\columnwidth]{figures/external_pilots/rank_cpso_bbob_stageC_pilot_v1_final_best_value.pdf}
  \caption{Ranghi finali appaiati nel pilot COCO BBOB; analisi descrittiva.}
  \label{fig:bbob-pilot-ranks}
\end{figure}}{\noindent\fbox{\parbox{0.95\columnwidth}{\centering Da generare: report COCO/cocopp e ranghi descrittivi BBOB.}}}

\subsection{Benchmark esterno: CEC 2013 LSGO}
\label{sec:results-cec}

Da riportare: andamento best-so-far ai checkpoint FE, statistiche per funzione,
mediane aggregate e confronti appaiati. I risultati del pilot a sei funzioni,
tre semi e 120\,000 FE saranno esplicitamente descritti come preliminari; le
conclusioni sulla suite richiederanno tutte le funzioni, i seed e il budget
confirmatory.

\IfFileExists{figures/external_pilots/rank_cpso_cec2013_stageC_pilot_v1_best_at_120000_FE.pdf}{%
\begin{figure}[t]
  \centering
  \includegraphics[width=\columnwidth]{figures/external_pilots/rank_cpso_cec2013_stageC_pilot_v1_best_at_120000_FE.pdf}
  \caption{Ranghi al checkpoint finale del pilot CEC 2013 LSGO.}
  \label{fig:cec-pilot-ranks}
\end{figure}}{\noindent\fbox{\parbox{0.95\columnwidth}{\centering Da generare: riepilogo CEC ai checkpoint FE.}}}

\subsection{Sintesi dei risultati}
\label{sec:results-summary}

Questa sottosezione conterrà una sintesi delle risposte alle RQ1--RQ5,
separando invarianti dimostrate e verificate, osservazioni empiriche e ipotesi
non supportate. Non saranno formulate graduatorie universali sulla base dei
soli pilot locali.
